%!Tex Program = xelatex
% -*-coding: utf-8 -*-
%  \author{初砚硕}

\def \xuewei{Doctor} % 定义学位 Doctor or Master

\def\xueke{Science} % 定义学科 Engineering, Science, Management, or Arts

\documentclass[cs4size,openany,twoside,UTF8,normalindentfirst,nofonts]{ctexbook}

% !Mode:: "TeX:UTF-8" 

\makeatletter
\@tempcnta=128
\loop \catcode\@tempcnta=13 \ifnum\@tempcnta<255 \advance \@tempcnta \@ne
\repeat
\makeatother

\newif\ifxueweidoctor %判断论文类型
\newif\ifxueweimaster
\def\temp{Doctor}
\ifx\temp\xuewei
  \xueweidoctortrue  \xueweimasterfalse
\fi
\def\temp{Master}
\ifx\temp\xuewei
  \xueweidoctorfalse  \xueweimastertrue
\fi

\ifxueweidoctor
  \newcommand{\cxuewei}{博士}
  \newcommand{\exuewei}{Doctor}
  \newcommand{\exueweier}{Doctoral}
  \newcommand{\xueweishort}{博}
\fi

\ifxueweimaster
  \newcommand{\cxuewei}{硕士}
  \newcommand{\exuewei}{Master}
  \newcommand{\exueweier}{Master's}
  \newcommand{\xueweishort}{硕}
\fi    % 硕博类型
%!Tex Program = xelatex
% -*-coding: utf-8 -*-


\usepackage{graphicx}

%\usepackage[a4paper,text={162true mm,249true mm},left=24true mm,
%            head=24true mm, top=24true mm, bottom=24true mm,
%           headheight=14true mm, headsep=10true mm
%           ]{geometry}% 东北林业大学要求的宽度

%\usepackage[a4paper,]{geometry}%四周都是24mm
\usepackage[a4paper,
left=24true mm,right=24true mm, top=25true mm, bottom=23true mm,
headheight=14true mm,
headsep=7true mm,
footskip=7true mm]{geometry}% 东北林业大学要求的宽度
\usepackage{layout}


\usepackage{titlesec}               % 控制标题的宏包
\usepackage{titletoc}                   % 控制目录的宏包
\usepackage{fancyhdr}                   % fancyhdr宏包 页眉和页脚的相关定义
%\usepackage[UTF8]{ctex}
\usepackage{color}          % 支持彩色
\usepackage{amsmath}        % AMSLaTeX宏包 用来排出更加漂亮的公式
\usepackage{amssymb}
\usepackage{epsfig}         % eps图像
\usepackage[below]{placeins}%允许上一个section 的浮动图形出现在下一个section的开始部分,还提供\FloatBarrier命令,使所有未处理的浮动图形立即被处理
\usepackage{flafter}       % 使得所有浮动体不能被放置在其浮动环境之前，以免浮动体在引述它的文本之前出现.
\usepackage{multirow}       %使用Multirow宏包，使得表格可以合并多个row 格
\usepackage{booktabs}       % 表格，横的粗线；\specialrule{1pt}{0pt}{0pt}
\usepackage{longtable}      %支持跨页的表格。
\usepackage{tabularx}
\usepackage{subfigure}%支持子图 %centerlast 设置最后一行是否居中
\usepackage[subfigure]{ccaption} %支持双语标题
\usepackage[sort&compress,numbers]{natbib}% 支持引用缩写的宏包
\makeatletter
%林大的参考文献研究用波浪链接~
\def\NAT@citexnum[#1][#2]#3{%
  \NAT@reset@parser
  \NAT@sort@cites{#3}%
  \NAT@reset@citea
  \@cite{\def\NAT@num{-1}\let\NAT@last@yr\relax\let\NAT@nm\@empty
    \@for\@citeb:=\NAT@cite@list\do
    {\@safe@activestrue
     \edef\@citeb{\expandafter\@firstofone\@citeb\@empty}%
     \@safe@activesfalse
     \@ifundefined{b@\@citeb\@extra@b@citeb}{%
       {\reset@font\bfseries?}
        \NAT@citeundefined\PackageWarning{natbib}%
       {Citation `\@citeb' on page \thepage \space undefined}}%
     {\let\NAT@last@num\NAT@num\let\NAT@last@nm\NAT@nm
      \NAT@parse{\@citeb}%
      \ifNAT@longnames\@ifundefined{bv@\@citeb\@extra@b@citeb}{%
        \let\NAT@name=\NAT@all@names
        \global\@namedef{bv@\@citeb\@extra@b@citeb}{}}{}%
      \fi
      \ifNAT@full\let\NAT@nm\NAT@all@names\else
        \let\NAT@nm\NAT@name\fi
      \ifNAT@swa
       \@ifnum{\NAT@ctype>\@ne}{%
        \@citea
        \NAT@hyper@{\@ifnum{\NAT@ctype=\tw@}{\NAT@test{\NAT@ctype}}{\NAT@alias}}%
       }{%
        \@ifnum{\NAT@cmprs>\z@}{%
         \NAT@ifcat@num\NAT@num
          {\let\NAT@nm=\NAT@num}%
          {\def\NAT@nm{-2}}%
         \NAT@ifcat@num\NAT@last@num
          {\@tempcnta=\NAT@last@num\relax}%
          {\@tempcnta\m@ne}%
         \@ifnum{\NAT@nm=\@tempcnta}{%
          \@ifnum{\NAT@merge>\@ne}{}{\NAT@last@yr@mbox}%
         }{%
           \advance\@tempcnta by\@ne
           \@ifnum{\NAT@nm=\@tempcnta}{%
             \ifx\NAT@last@yr\relax
               \def@NAT@last@yr{\@citea}%
             \else
               \def@NAT@last@yr{$\sim$\NAT@penalty}%
             \fi
           }{%
             \NAT@last@yr@mbox
           }%
         }%
        }{%
         \@tempswatrue
         \@ifnum{\NAT@merge>\@ne}{\@ifnum{\NAT@last@num=\NAT@num\relax}{\@tempswafalse}{}}{}%
         \if@tempswa\NAT@citea@mbox\fi
        }%
       }%
       \NAT@def@citea
      \else
        \ifcase\NAT@ctype
          \ifx\NAT@last@nm\NAT@nm \NAT@yrsep\NAT@penalty\NAT@space\else
            \@citea \NAT@test{\@ne}\NAT@spacechar\NAT@mbox{\NAT@super@kern\NAT@@open}%
          \fi
          \if*#1*\else#1\NAT@spacechar\fi
          \NAT@mbox{\NAT@hyper@{{\citenumfont{\NAT@num}}}}%
          \NAT@def@citea@box
        \or
          \NAT@hyper@citea@space{\NAT@test{\NAT@ctype}}%
        \or
          \NAT@hyper@citea@space{\NAT@test{\NAT@ctype}}%
        \or
          \NAT@hyper@citea@space\NAT@alias
        \fi
      \fi
     }%
    }%
      \@ifnum{\NAT@cmprs>\z@}{\NAT@last@yr}{}%
      \ifNAT@swa\else
        \@ifnum{\NAT@ctype=\z@}{%
          \if*#2*\else\NAT@cmt#2\fi
        }{}%
        \NAT@mbox{\NAT@@close}%
      \fi
  }{#1}{#2}%
}%
\makeatother

\usepackage{enumitem}       %使用enumitem宏包,改变列表项的格式
\usepackage{calc}           %长度可以用+ - * / 进行计算


%%\usepackage{txfonts}
\usepackage{mathtools}

\usepackage{bm}              % 处理数学公式中的黑斜体的宏包

\usepackage[xetex,
            bookmarksnumbered=true,
            bookmarksopen=true,
            colorlinks=false,
            pdfborder={0 0 1},
            citecolor=blue,
            linkcolor=red,
            anchorcolor=green,
            urlcolor=blue,
            breaklinks=true,
            naturalnames  %与algorithm2e宏包协调
            ]{hyperref}
%\usepackage{fontenc}
%xunicode xunicode Ross Moore’s xunicode package is now automatically loaded for users of both XELATEX and LuaLATEX.
\usepackage[BoldFont,SlantFont]{xeCJK}
\defaultfontfeatures{Mapping=tex-text}

\xeCJKsetemboldenfactor{1}%只对随后定义的CJK字体有效
\setCJKfamilyfont{hei}{SimHei}
\xeCJKsetemboldenfactor{4}
\setCJKfamilyfont{song}{SimSun}
\xeCJKsetemboldenfactor{1}
\setCJKfamilyfont{fs}{FangSong}
\setCJKfamilyfont{kai}{KaiTi}
\setCJKfamilyfont{li}{LiSu}
\setCJKfamilyfont{xw}{STXinwei}

\setCJKmainfont{SimSun}
\setmainfont{Times New Roman}
\setsansfont{Arial}

\newcommand{\hei}{\CJKfamily{hei}}% 黑体   (Windows自带simhei.ttf)
\newcommand{\song}{\CJKfamily{song}}    % 宋体   (Windows自带simsun.ttf)
\newcommand{\fs}{\CJKfamily{fs}}        % 仿宋体 (Windows自带simfs.ttf)
\newcommand{\kai}{\CJKfamily{kai}}      % 楷体   (Windows自带simkai.ttf)
\newcommand{\li}{\CJKfamily{li}}        % 隶书   (Windows自带simli.ttf)
\newcommand{\xw}{\CJKfamily{xw}}        % 隶书   (Windows自带simli.ttf)


\newfontfamily\arial{Arial}
\newfontfamily\timesnewroman{Times New Roman}

\usepackage[boxed,linesnumbered,algochapter]{algorithm2e}%\let\chapter\undefined
% 算法的宏包，注意宏包兼容性，先后顺序为float、hyperref、algorithm(2e)，否则无法生成算法列表

\usepackage{setspace}%更改行距的宏包
%\usepackage{syntonly}检查语法不编译
%\syntaxonly
\usepackage{tikz}%R语言作图使用
\usepackage{xltxtra}%输出XeLaTeX
 % 引用的宏包

\graphicspath{{figures/}} %定义所有的eps文件在 figures 子目录下


\begin{document}

%!Tex Program = xelatex
% -*-coding: utf-8 -*-


\newcommand{\xiaochu}{\fontsize{36pt}{36pt}\selectfont}       % 一号, 1.倍行距
\newcommand{\yihao}{\fontsize{26pt}{26pt}\selectfont}       % 一号, 1.倍行距
\newcommand{\xiaoyi}{\fontsize{24pt}{24pt}\selectfont}      % 小一, 1.倍行距
\newcommand{\erhao}{\fontsize{22pt}{22pt}\selectfont}       % 二号, 1.倍行距
\newcommand{\xiaoer}{\fontsize{18pt}{18pt}\selectfont}      % 小二, 单倍行距
\newcommand{\sanhao}{\fontsize{16pt}{16pt}\selectfont}      % 三号, 1.倍行距
\newcommand{\xiaosan}{\fontsize{15pt}{15pt}\selectfont}     % 小三, 1.倍行距
\newcommand{\sihao}{\fontsize{14pt}{14pt}\selectfont}       % 四号, 1.0倍行距
\newcommand{\xiaosi}{\fontsize{12pt}{12pt}\selectfont}      % 小四, 1.倍行距
\newcommand{\wuhao}{\fontsize{10.5pt}{10.5pt}\selectfont}   % 五号, 单倍行距
\newcommand{\xiaowu}{\fontsize{9pt}{9pt}\selectfont}        % 小五, 单倍行距



%避免宏包 hyperref 和 arydshln 不兼容带来的目录链接失效的问题。
\def\temp{\relax}
\let\temp\addcontentsline
\gdef\addcontentsline{\phantomsection\temp}

\makeatletter
\gdef\hitempty{}%全局定义空，貌似没用

%重新定义BiChapter命令，可实现标题手动换行，但不影响目录
\def\BiChapter{\relax\@ifnextchar [{\@BiChapter}{\@@BiChapter}}
\def\@BiChapter[#1]#2#3{\chapter[#1]{#2}
    \addcontentsline{toe}{chapter}{\bfseries \xiaosi Chapter \thechapter\hspace{0.5em} #3}}
\def\@@BiChapter#1#2{\chapter{#1}
    \addcontentsline{toe}{chapter}{\bfseries \xiaosi Chapter \thechapter\hspace{0.5em}{\boldmath #2}}}

\newcommand{\BiSection}[2]
{   \section{#1}
    \addcontentsline{toe}{section}{\protect\numberline{\csname thesection\endcsname}#2}
}

\newcommand{\BiSubsection}[2]
{    \subsection{#1}
    \addcontentsline{toe}{subsection}{\protect\numberline{\csname thesubsection\endcsname}#2}
}

\newcommand{\BiSubsubsection}[2]
{    \subsubsection{#1}
    \addcontentsline{toe}{subsubsection}{\protect\numberline{\csname thesubsubsection\endcsname}#2}
}

\newcommand{\BiAppendixChapter}[2] % 该附录命令适用于发表文章，简历等
{\phantomsection
\markboth{#1}{#1}
\addcontentsline{toc}{chapter}{\xiaosi #1}
\addcontentsline{toe}{chapter}{\bfseries \xiaosi #2}  \chapter*{#1}
}

\newcommand{\BiAppChapter}[2]    % 该附录命令适用于有章节的完整附录
{\phantomsection  \chapter{#1}   %\markboth{\MakeUppercase{#1}}{\MakeUppercase{#1}} %为了winedt中project tree 中toc正确显示，不要挪到下一行；
\addcontentsline{toe}{chapter}{\bfseries \xiaosi Appendix \thechapter~~#2}
}

\renewcommand{\thefigure}{\arabic{chapter}-\arabic{figure}}%使图编号为 7-1 的格式 %\protect{~}
\renewcommand{\thesubfigure}{\alph{subfigure})}%使子图编号为 a)的格式
\renewcommand{\p@subfigure}{\thefigure~} %%使子图引用为 7-1 a) 的格式，母图编号和子图编号之间用~加一个空格
\renewcommand{\thetable}{\arabic{chapter}-\arabic{table}}%%使表编号为 7-1 的格式
\renewcommand{\theequation}{\arabic{chapter}-\arabic{equation}}%% 使公式编号为 7-1 的格式

\newcommand{\algoenname}{Algo.} %算法英文标题
\newfloatlist[chapter]{algoen}{aen}{\listalgoenname}{\algoenname}
\newfixedcaption{\algoencaption}{algoen}
\renewcommand{\thealgoen}{\thechapter-\arabic{algocf}}
\renewcommand{\@cftmakeaentitle}{\chapter*{\listalgoenname\@mkboth{\bfseries\listalgoenname}{\bfseries\listalgoenname}}}

\renewcommand{\algorithmcfname}{算法}
\setlength\AlCapSkip{1.2ex}
\SetAlgoSkip{1pt}
\renewcommand{\algocf@captiontext}[2]{\wuhao#1\algocf@typo ~ \AlCapFnt{}#2} % text of caption
\expandafter\ifx\csname algocf@within\endcsname\relax% if \algocf@within doesn't exist
\renewcommand\thealgocf{\@arabic\c@algocf} % and the way it is printed
\else%                                    else
\renewcommand\thealgocf{\csname the\algocf@within\endcsname-\@arabic\c@algocf}
\fi
\renewcommand{\algocf@makecaption}[2]{%中英文双标题一定多于一行，因此去掉单行时的判断，并将\parbox中标题设置为居中
  \addtolength{\hsize}{\algomargin}%
  \sbox\@tempboxa{\algocf@captiontext{#1}{#2}}%
    \hskip .5\algomargin%
    \parbox[t]{\hsize}{\centering\algocf@captiontext{#1}{#2}}% 
  \addtolength{\hsize}{-\algomargin}%
}
\newcommand{\AlgoBiCaption}[2]{%直接取出自定义的中英文标题条目加入到真正的\caption 中  
   \caption[#1]{\protect\setlength{\baselineskip}{1.5em}#1 \protect \\ Algo. \thealgocf~~ #2} % \algoencaption{#2}   
   \addcontentsline{aen}{algoen}{\protect\numberline{\thealgoen}{#2}}
   }
\makeatother

%定义 学科 学位
\def \xuekeEngineering {Engineering}
\def \xuekeScience {Science}
\def \xuekeManagement {Management}
\def \xuekeArts {Arts}

\ifx \xueke \xuekeEngineering
\newcommand{\cxueke}{工学}
\newcommand{\exueke}{Engineering}
\fi

\ifx \xueke \xuekeScience
\newcommand{\cxueke}{理学}
\newcommand{\exueke}{Science}
\fi

\ifx \xueke \xuekeManagement
\newcommand{\cxueke}{管理学}
\newcommand{\exueke}{Management}
\fi

\ifx \xueke \xuekeArts
\newcommand{\cxueke}{文学}
\newcommand{\exueke}{Arts}
\fi

\ifx \xueke \xuekePhilosophy
\newcommand{\cxueke}{哲学}
\newcommand{\exueke}{Philosophy}
\fi 

\ifx \xueke \xuekeEconomics
\newcommand{\cxueke}{经济学}
\newcommand{\exueke}{Economics}
\fi 

\ifx \xueke \xuekeLaws
\newcommand{\cxueke}{法学}
\newcommand{\exueke}{Laws}
\fi 

\ifx \xueke \xuekeEducation
\newcommand{\cxueke}{教育学}
\newcommand{\exueke}{Education}
\fi 

\ifx \xueke \xuekeHistory
\newcommand{\cxueke}{历史学}
\newcommand{\exueke}{History}
\fi 
 % 文本格式定义
%!Tex Program = xelatex
% -*-coding: utf-8 -*-

\theoremstyle{plain}
\theorembodyfont{\song\rmfamily}
\theoremheaderfont{\hei\rmfamily}
\newtheorem{definition}{\hei 定义}[chapter]
\newtheorem{example}{\hei 例}[chapter]
\newtheorem{algo}{\hei 算法}[chapter]
\newtheorem{theorem}{\hei 定理}[chapter]
\newtheorem{axiom}{\hei 公理}[chapter]
\newtheorem{proposition}{\hei 命题}[chapter]
\newtheorem{lemma}{\hei 引理}[chapter]
\newtheorem{corollary}{\hei 推论}[chapter]
\newtheorem{remark}{\hei 注解}[chapter]
\newenvironment{proof}{\noindent{\hei 证明：}}{\hfill $ \square $ \vskip 4mm}
\theoremsymbol{$\square$}
\setlength{\theorempreskipamount}{0pt}
\setlength{\theorempostskipamount}{-2pt}

\allowdisplaybreaks[4]

\setlength{\parindent}{2em}

\renewcommand\contentsname{目录}%目~~~~录

\CTEXsetup[number={\arabic{chapter}}]{chapter}
\renewcommand\chaptername{~\thechapter~}

\setcounter{secnumdepth}{4} \setcounter{tocdepth}{2}


\titleformat{\chapter}{\center\xiaoer\bf\hei}{\chaptername}{0.5em}{}
\titlespacing{\chapter}{0pt}{-6.35mm}{15pt}
%\titleformat{\chapter*}{\center\xiaoer\bf\hei\arial}{\chaptername}{0.5em}{}
%\titlespacing{\chapter*}{0pt}{-6.35mm}{15pt}

\titleformat{\section}{\xiaosan\bf\hei\arial}{\thesection}{0.5em}{}
\titlespacing{\section}{0pt}{6pt}{6pt}
\titleformat{\subsection}{\sihao\bf\hei\arial}{\thesubsection}{0.5em}{}
\titlespacing{\subsection}{0pt}{6pt}{0mm}
\titleformat{\subsubsection}[runin]{\song\arial}{\thesubsubsection}{0.5em}{}[\quad\quad ]
\titlespacing{\subsubsection}{0pt}{0pt}{0pt}

\titlecontents{chapter}[0em]{\bf\hei}{\arial\thecontentslabel~}{}{\bf\titlerule*[4pt]{.}\contentspage}
\titlecontents{section}[0em]{}{\thecontentslabel~}{}{\titlerule*[4pt]{.}\contentspage}
\titlecontents{subsection}[0em]{}{\thecontentslabel~}{}{\titlerule*[4pt]{.}\contentspage}

%\dottedcontents{section}[40pt]{}{22pt}{0.3pc}
%\dottedcontents{subsection}[62pt]{}{32pt}{0.3pc}



% 缩小目录中各级标题之间的缩进，使它们相隔一个字符距离，也就是12pt
\makeatletter
\renewcommand*\l@chapter{\@dottedtocline{0}{0em}{4em}}% 控制英文目录： 细点\@dottedtocline  粗点\@dottedtoclinebold
\renewcommand*\l@section{\@dottedtocline{1}{1.5em}{1.8em}}
\renewcommand*\l@subsection{\@dottedtocline{2}{2.5em}{2.5em}}


% 定义页眉和页脚 使用fancyhdr 宏包
%\newcommand{\makeheadrule}{
%    \rule[0pt]{\textwidth}{0.75pt}\\[-16.5pt]
%\rule{\textwidth}{0.8mm}
%}
%\renewcommand{\headrule}{
%    {\if@fancyplain\let\headrulewidth\plainheadrulewidth\fi
%     \makeheadrule}
%}
%\renewcommand{\headrule}{
%    \hrule height0.4true mm width\headwidth \vskip0.2true mm
%    \hrule height0.2true mm width\headwidth \vskip-0.8true mm
%}

\renewcommand{\headrule}{
    \vskip0.5true mm
    \hrule height0.5true mm width\textwidth \vskip0.3true mm
    \hrule height0.3true mm width\textwidth \vskip-1.6true mm
}

\pagestyle{fancyplain}

%去掉章节标题中的数字
%%不要注销这一行，否则页眉会变成：“第1章1  绪论”样式
\renewcommand{\chaptermark}[1]{\markboth{\chaptertitlename~\ #1}{}}
\fancyhf{}

%在book文件类别下,\leftmark自动存录各章之章名,\rightmark记录节标题
%% 页眉字号 要求 小五
%根据单双面打印设置不同的页眉；
%\fancyhead[CO]{\song \xiaowu 东北林业大学硕士学位论文}
%\fancyhead[CE]{\song \xiaowu \leftmark}%
\fancyfoot[C,C]{\xiaowu -~\thepage~-}


\renewcommand\frontmatter{\cleardoublepage
  \@mainmatterfalse
  \pagenumbering{Roman}}

% 设置行距和段落间垂直距离
\renewcommand{\CJKglue}{\hskip 0.56pt plus 0.08\baselineskip} % 加大字间距，使每行34个字

% 调整列表环境的垂直间距
\setitemize{itemindent=3em,leftmargin=0pt,itemsep=0ex,listparindent=2em,partopsep=0pt,parsep=0ex,topsep=0ex}
\setenumerate{itemindent=3em,leftmargin=0pt,itemsep=0ex,listparindent=2em,partopsep=0pt,parsep=0ex,topsep=0ex}
\setdescription{itemindent=3em,leftmargin=0pt,itemsep=0ex,listparindent=2em,partopsep=0pt,parsep=0ex,topsep=0ex}


\renewcommand\@biblabel[1]{[#1]\hspace{0.5em}} % 不去除参考文献里标号两边的括号
\newcommand{\ucite}[1]{$^{\mbox{\scriptsize \cite{#1}}}$} % 增加 \ucite 命令使显示的引用为上标形式
\newcommand{\citeup}[1]{$^{\mbox{\scriptsize \cite{#1}}}$} % for WinEdt users

% 定制浮动图形和表格标题样式
\renewcommand{\figurename}{图}
\renewcommand{\tablename}{表}
\captionnamefont{\wuhao}
\captiontitlefont{\wuhao}
\captiondelim{~~}
\renewcommand{\subcapsize}{\wuhao}
\setlength{\abovecaptionskip}{0pt}
\setlength{\belowcaptionskip}{0pt}
%将章标题中的中文数字（一、二、三）变为阿拉伯数字(1,2,3)

%%\renewcommand\CJKthechapter{{\@arabic\c@chapter}}

% 自定义项目列表标签及格式 \begin{nefulist} 列表项 \end{nefulist}
\newcounter{nefuctr} %自定义新计数器
\newenvironment{nefulist}{%%%%%定义新环境
\begin{list}{{(\arabic{nefuctr})}} %%标签格式
    {
     \usecounter{nefuctr}
     \setlength{\leftmargin}{0cm}     %左边界
     \setlength{\parsep}{0ex}         %段落间距
     \setlength{\topsep}{0pt}         %列表到上下文的垂直距离
     \setlength{\itemsep}{0ex}        %标签间距
     \setlength{\labelwidth}{0.5em}     %标号和列表项之间的距离,默认0.5em
     \setlength{\labelsep}{0.5em}     %标号和列表项之间的距离,默认0.5em
     \setlength{\itemindent}{3em}    % 标签缩进量
     \setlength{\listparindent}{2em} % 段落缩进量
    }}
{\end{list}}%%%%%

% 自定义项目列表标签及格式 \begin{publist} 列表项 \end{publist}
\newcounter{pubctr} %自定义新计数器
\newenvironment{publist}{%%%%%定义新环境
\begin{list}{\arabic{pubctr}} %%标签格式
    {
     \usecounter{pubctr}
     \setlength{\leftmargin}{1.5em}     % 左边界 \leftmargin =\itemindent + \labelwidth + \labelsep
     \setlength{\itemindent}{0em}     % 标号缩进量
     \setlength{\labelsep}{1em}       % 标号和列表项之间的距离,默认0.5em
     \setlength{\rightmargin}{0em}    % 右边界
     \setlength{\topsep}{0ex}         % 列表到上下文的垂直距离
     \setlength{\parsep}{0ex}         % 段落间距
     \setlength{\itemsep}{0ex}        % 标签间距
     \setlength{\listparindent}{0pt} % 段落缩进量
    }}
{\end{list}}%%%%%

% 默认字体
\renewcommand\normalsize{
  \@setfontsize\normalsize{12pt}{12pt}
  \setlength\abovedisplayskip{4pt}
  \setlength\abovedisplayshortskip{4pt}
  \setlength\belowdisplayskip{\abovedisplayskip}
  \setlength\belowdisplayshortskip{\abovedisplayshortskip}
  \let\@listi\@listI}
\def\defaultfont{\renewcommand{\baselinestretch}{1.57}\normalsize\selectfont}
\predisplaypenalty=0  %公式之前可以换页，公式出现在页面顶部

% 封面、摘要、版权、致谢格式定义
\def\ctitle#1{\def\@ctitle{#1}}\def\@ctitle{}
\def\cdegree#1{\def\@cdegree{#1}}\def\@cdegree{}
\def\caffil#1{\def\@caffil{#1}}\def\@caffil{}
\def\csubject#1{\def\@csubject{#1}}\def\@csubject{}
\def\cauthor#1{\def\@cauthor{#1}}\def\@cauthor{}
\def\csupervisor#1{\def\@csupervisor{#1}}\def\@csupervisor{}
\def\cchief#1{\def\@cchief{#1}}\def\@cchief{}
\def\creview#1{\def\@creview{#1}}\def\@creview{}


\def\cassosupervisor#1{\def\@cassosupervisor{~ & {\hei 副 \hfill 导 \hfill 师：} & #1\\}}\def\@cassosupervisor{}
\def\ccosupervisor#1{\def\@ccosupervisor{~ & {\hei 联 \hfill 合\hfill 导 \hfill 师：} & #1\\}}\def\@ccosupervisor{}
\def\cdate#1{\def\@cdate{#1}}\def\@cdate{}% 论文提交日期
\def\cdatequestion#1{\def\@cdatequestion{#1}}\def\@cdatequestion{}% 答辩日期
\def\cdatedegree#1{\def\@cdatedegree{#1}}\def\@cdatedegree{}% 授予学位日期

\long\def\cabstract#1{\long\def\@cabstract{#1}}\long\def\@cabstract{}
\def\ckeywords#1{\def\@ckeywords{#1}}\def\@ckeywords{}

\def\etitle#1{\def\@etitle{#1}}\def\@etitle{}
\def\edegree#1{\def\@edegree{#1}}\def\@edegree{}
\def\eaffil#1{\def\@eaffil{#1}}\def\@eaffil{}
\def\esubject#1{\def\@esubject{#1}}\def\@esubject{}
\def\eauthor#1{\def\@eauthor{#1}}\def\@eauthor{}
\def\esupervisor#1{\def\@esupervisor{#1}}\def\@esupervisor{}
%\def\eassosupervisor#1{\def\@eassosupervisor{#1}}\def\@eassosupervisor{}
\def\eassosupervisor#1{\def\@eassosupervisor{~ & \textbf{Associate Supervisor:} & #1\\}}\def\@eassosupervisor{}
%\def\ecosupervisor#1{\def\@ecosupervisor{#1}}\def\@ecosupervisor{}
\def\ecosupervisor#1{\def\@ecosupervisor{~ & \textbf{Co Supervisor:} & #1\\}}\def\@ecosupervisor{}
\def\edate#1{\def\@edate{#1}}\def\@edate{}
\def\edatequestion#1{\def\@edatequestion{#1}}\def\@edatequestion{}

\long\def\eabstract#1{\long\def\@eabstract{#1}}\long\def\@eabstract{}
\long\def\NotationList#1{\long\def\@NotationList{#1}}\long\def\@NotationList{}
\def\ekeywords#1{\def\@ekeywords{#1}}\def\@ekeywords{}
\def\natclassifiedindex#1{\def\@natclassifiedindex{#1}}\def\@natclassifiedindex{}
\def\internatclassifiedindex#1{\def\@internatclassifiedindex{#1}}\def\@internatclassifiedindex{}
\def\statesecrets#1{\def\@statesecrets{#1}}\def\@statesecrets{}

% 定义封面
\def\makecover{
      \begin{titlepage}
            % 封面一
            %\vspace*{0.3cm}
%            \begin{center}{
%                  \xiaoyi
%                  \begin{center}{
%                        \song\bfseries 硕士学位论文
%                  }\end{center}
%            }
%
%            \vspace{0.6cm}
%
%            \parbox[t][2.9cm][t]{\textwidth}{
%                  \erhao
%                  \begin{center}{
%                        \hei\@ctitle
%                  }\end{center}
%            }
%
%            \parbox[t][5.1cm][t]{\textwidth}{
%                  \erhao %英文标题太长时可以采用\xiaoer
%                  \begin{center}{
%                        \bfseries\@etitle
%                  }\end{center}
%            }
%
%            \parbox[t][8.2cm][t]{\textwidth}{
%                  \xiaoer
%                  \begin{center}{
%                        \song\bfseries  \@cauthor
%                  }\end{center}
%            }
%
%            \parbox[t][1cm][t]{\textwidth}{
%                  \begin{center}{
%                        \song \xiaoer\bfseries \@cdate
%                  }\end{center}
%            }\end{center}
%%%%%%%%%%%%%%%%%%%%%%%%%%%%%%%%%%%%%%%%%%%%%%%%%%%%%%%%%%%%%%%%%%%%%%%%%%%%%%%%%%%%%%%%%%%%%%%%%%%%%%%%%%%%%%%%%%%%%%%%%
            %内封
            \newpage
            \thispagestyle{empty}
                        \newlength{\widthtemp}
                        \settowidth{\widthtemp}{\song \sihao 学校代码：}
                        {\setlength{\parindent}{2.2em}
                        \parbox[t][17mm][t]{\textwidth}
                        {\begin{spacing}{1.25}
                              \hspace*{10.34cm}{\song \sihao 学校代码：10213}\song \xiaoer\\
                              \indent\hspace*{10.34cm}\makebox[\widthtemp][s]{\song \sihao 学\hfill 号：}\song \sihao  05001\song \xiaoer\\
                        \end{spacing}}}
\begin{center}
            \vspace{3.25mm}
            \parbox[t][3.3cm][t]{\textwidth}{\begin{center}{\xiaochu\xw 学\ \ 位\ \ 论\ \ 文}\end{center}}
            \parbox[t][3.4cm][t]{\textwidth}{\begin{center}\erhao\hei\@ctitle\end{center}}
            \parbox[t][3.75cm][t]{\textwidth}{\begin{center}\xiaoyi\kai\@cauthor\end{center}}
            \newlength{\lsspace}
            \setlength{\lsspace}{0.78cm}
      \begin{minipage}[t]{\textwidth}
            {\sihao%\hspace{0.56cm}
                              \song %\renewcommand{\arraystretch}{0.6}
                              \settowidth{\widthtemp}{\hei \sihao 授予学位日期}
                              \vspace{0.6em}\hspace{1ex}
                              \renewcommand{\arraystretch}{0}
\begin{tabular}{r@{：}lr@{：}l}
{\hei \rule{0pt}{\lsspace}指导教师姓名}&\@csupervisor&\multicolumn{2}{l}{东北林业大学}\\
{\hei \rule{0pt}{\lsspace}申请学位级别}&\@cdegree &\makebox[\widthtemp][s]{\hei 学\hfill 科\hfill 专\hfill 业}&\@csubject\\
{\hei \rule{0pt}{\lsspace}论文提交日期}&\@cdate  &{\hei 论文答辩日期}&\@cdatequestion  \\% \@cdegree\\
{\hei \rule{0pt}{\lsspace}授予学位单位}&东北林业大学&{\hei 授予学位日期}&\@cdatedegree \\
\multicolumn{4}{l}{\rule{0pt}{9mm}}
\end{tabular}}
      \end{minipage}
      {
      \settowidth{\widthtemp}{\hei\sihao 答辩委员会主席}
                                  \sihao\song
                                  \hspace*{5em}\begin{tabular}{c@{：}l}
                                    {\hei\sihao 答辩委员会主席} & {\@cchief}\\
                                    {\rule{0pt}{0.8\lsspace}\makebox[\widthtemp][s]{\hei\sihao 论\hfill\ 文\hfill\ 评\hfill\ 阅\hfill\ 人}} & {\@creview}\\
                                  \end{tabular}
      }
      \vspace{14mm}
                         \begin{center}
                     　　\includegraphics[height=1.53cm,width=7.00cm,bb=0 0 887 191]{nefupic/nefu_pic.png}
                  　　   \end{center}
\end{center}
%\@cauthor
%%%%%%%%%%%%%%%%%%%%%%%%%%%%%%%%%%%%%%%%%%%%%%%%%%%%%%%%%%%%%%%%%%%%%%%%%%%%%%%%%%%%%%%%%%%%%%%%%%%%%%%%%%%%%%%%%%%%%%%%%
            % 英文封面
            \newpage
            \thispagestyle{empty}
            \vspace*{1.5mm}
            {\setlength{\parindent}{23.62em}
                        \parbox[t][29mm][t]{\textwidth}
                        {\begin{spacing}{1.3}
                              {\xiaosi University Code：10213}\song \xiaosan\\
                              {\indent\xiaosi Register Code ：05001}
                        \end{spacing}}}
            \begin{center}
                  \parbox[t][41mm][t]{\textwidth}{\xiaoer
                        \begin{center} { Dissertation for the Degree of Master}\end{center}
                  } %与中文保持一致，删除in {\exueke}
                  \parbox[t][68mm][t]{\textwidth}{\erhao
                        \begin{center} {
                              \@etitle
                        }\end{center}
                  }
                  {
                        {
                              \sihao\noindent\hspace*{6.6mm}
                              \renewcommand{\arraystretch}{1.5}
                                    \begin{tabularx}{\textwidth}{@{}l@{~}X@{}}
                                          \textbf{Candidate:}                     &  \@eauthor\\
                                          \textbf{Supervisor:}                    &  \@esupervisor\\
                                          \textbf{Associate Supervisor:}          &  \@eassosupervisor\\
                                          \textbf{Academic Degree Applied for:}   &  \@edegree\\
                                          \textbf{Specialty:}                     &  \@esubject\\
                                          \textbf{Date of Oral Examination:}      &  \@edatequestion\\
                                          \textbf{University:} &  Northeast Forestry University
                                    \end{tabularx}
                              \renewcommand{\arraystretch}{1}

                        }
                  }
            \end{center}
      \end{titlepage}

      %%%%%%%%%%%%%%%%%%%   Abstract and keywords  %%%%%%%%%%%%%%%%%%%%%%%
      \clearpage
      \chapter*{摘要}
      \phantomsection
      \markboth{摘要}{摘要}
      %\vspace{1em}
      \addcontentsline{toc}{chapter}{\xiaosi 摘要}
      \addcontentsline{toe}{chapter}{\bfseries \xiaosi Abstract (In Chinese)}
      %\BiAppendixChapter{摘要}{}
      \setcounter{page}{1}
      \song\defaultfont
      \@cabstract
      \vspace{\baselineskip}

      \hangafter=1\hangindent=52.3pt\noindent
      {\hei 关键词：} \@ckeywords

      %%%%%%%%%%%%%%%%%%%   English Abstract  %%%%%%%%%%%%%%%%%%%%%%%%%%%%%%
      \clearpage
      \BiAppendixChapter{\textbf\timesnewroman{Abstract}}{Abstract (In English)} % 不要挪到下一行，生成正确的摘要toc
      \@eabstract
      \vspace{\baselineskip}

      \hangafter=1\hangindent=60pt\noindent
      {\bfseries{Keywords:}}  \@ekeywords

}

%%%%%%%%%%%%%%%%%%%%%%%%%%%%%%%%%%%%%%%%%%%%%%%%%%%%%%%%%%%%%%%
% 英文目录格式
\def\@dotsep{0.75}           % 定义英文目录的点间距
\setlength\leftmargini {0pt}
\setlength\leftmarginii {0pt}
\setlength\leftmarginiii {0pt}
\setlength\leftmarginiv {0pt}
\setlength\leftmarginv {0pt}
\setlength\leftmarginvi {0pt}

\def\engcontentsname{\bfseries Contents}
\newcommand\tableofengcontents{
   \pdfbookmark[0]{Contents}{econtent}
     \@restonecolfalse
   \chapter*{\engcontentsname  %chapter*上移一行，避免在toc中出现。
       \@mkboth{%
          \engcontentsname}{\engcontentsname}}
   \@starttoc{toe}%
   \if@restonecol\twocolumn\fi
}

\urlstyle{same}  %论文中引用的网址的字体默认与正文中字体不一致，这里修正为一致的。

\renewcommand\endtable{\vspace{-4mm}\end@float}

\makeatother

\fancyhead[CO]{\song \xiaowu \leftmark}
\fancyhead[CE]{\song \xiaowu \leftmark}
\frontmatter
% -*-coding: utf-8 -*-

\newcommand{\chinesethesistitle}{东北林业大学硕/博士毕业论文模版} %授权书用，无需断行
\newcommand{\englishthesistitle}{\uppercase{research on some subject}} %\uppercase 作用：将英文标题字母全部大写；
\newcommand{\chinesethesistime}{2012/12}  %论文提交日期
\newcommand{\chinesethesistimequestion}{2013/5}  %答辩日期
\newcommand{\englishthesistimequestion}{May 20th, 2013}  %答辩日期

\newcommand{\cattaindegree}{2013/5/30}  %获得学位日期

\newcommand{\englishthesistime}{December, 2010}    %封面底部的日期英文形式

\ctitle{东林硕/博士毕业论文\XeLaTeX 模版}  %封面用论文标题，自己可手动断行
\cdegree{硕士}
\csubject{某某专业}                 %(~按二级学科填写~)
\caffil{某某学院} %（在校生填所在系名称，同等学力人员填工作单位）
\cauthor{某某某}
\csupervisor{某某某} %导师名字
\cchief{}
\creview{}

\cdate{\chinesethesistime}
\cdatequestion{\chinesethesistimequestion}
\cdatedegree{\cattaindegree}

\etitle{\englishthesistitle}
\edegree{Master of \exueke}
\esubject{Computer Science and Technology}  %英文二级学科名
\eaffil{College of  Information and Computer Engineering}%英文单位 %换行用\newline，不要用\\
\eauthor{ABCD}                   %作者姓名 （英文）
\esupervisor{Prof. EFGH}       % 导师姓名 （英文）
%\ecosupervisor{Professor X}
%\eassosupervisor{Professor Y}
\edate{\englishthesistime}
\edatequestion{\englishthesistimequestion}

\natclassifiedindex{TM301.2}  %国内图书分类号
\internatclassifiedindex{62-5}  %国际图书分类号
\statesecrets{公开} %秘密

\iffalse
\BiAppendixChapter{摘~~~~要}{}  %使用winedt编辑时文档结构图（toc）中为了显示摘要，故增加此句；
\fi
\cabstract{
东北林业大学硕博毕业论文\XeLaTeX 模版

}

\ckeywords{关键词~1；关键词~2；关键词~3；……；关键词~6（关键词总共~3~—~6~个，最后一个关键词后面没有标点符号）}

\eabstract{
Externally pressurized gas bearing has been widely used in the field of aviation, semiconductor, weave, and measurement apparatus because of its advantage of high accuracy, little friction, low heat distortion, long life-span, and no pollution. In this thesis, based on the domestic and overseas researching……

}

\ekeywords{keyword 1, keyword 2, keyword 3, ……, keyword 6 (no punctuation at the end)}

\makecover
\clearpage
 % 封面

%% 目录

\defaultfont
\clearpage{\pagestyle{empty}\cleardoublepage}
\pdfbookmark[0]{目~~~~录}{mulu}
\tableofcontents    % 中文目录
\clearpage{\pagestyle{empty}\cleardoublepage}

%\addtocontents{toc}{\protect\vskip1\baselineskip} % 中文目录增加空行

\ifxueweidoctor
\tableofengcontents % 英文目录
\clearpage{\pagestyle{empty}\cleardoublepage}     % 清除目录后面空页的页眉和页脚

\mainmatter\defaultfont\sloppy\raggedbottom
\fancyhead[CO]{\song \xiaowu 东北林业大学硕士学位论文}
\fancyhead[CE]{\song \xiaowu \leftmark}%
\titleformat{\chapter}{\center\xiaoer\bf\hei\arial}{\chaptername}{0.5em}{}% 实现英文摘要字体和正文章标题字体不一样
\titlespacing{\chapter}{0pt}{-6.35mm}{15pt}


%!Tex Program = xelatex
% -*-coding: utf-8 -*-
        \BiChapter{模版简介}{Review}

        \BiSection{模版来源}{Original of this template} 

本模版是从哈工大\LaTeX 模版2.0改版而来。

\BiSection{模版使用方法}{How to use this template} 
本模版的插图、公式、图表命令与哈工大\LaTeX 模版2.0一致。
其具体方式可查看哈工大\LaTeX 模版2.0的模版说明。
不同之处有如下几点：
\begin{nefulist}
  \item 列表命令
  
  \begin{verbatim}
    \begin{nefulist}
      \item *
      \item *
      \item *      
    \end{nefulist}
  \end{verbatim}
  \item 公式引用命令
  
  \begin{verbatim}
    \nefuEqRef{label}
  \end{verbatim}
\end{nefulist}


%!Tex Program = xelatex
% -*-coding: utf-8 -*-
\phantomsection
\markboth{结论}{结论}
\addcontentsline{toc}{chapter}{\xiaosi 结论}
\addcontentsline{toe}{chapter}{\bfseries \xiaosi Conclusions}  \chapter*{结论}

















   % 结论
%参考文献

\fancyhead[CO]{\song \xiaowu \leftmark}
\fancyhead[CE]{\song \xiaowu \leftmark}%
\defaultfont

\bibliographystyle{chinesebst2005_UTF8}%
\addcontentsline{toc}{chapter}{参考文献}      % 参考文献加入到中文目录
\addcontentsline{toe}{chapter}{\bfseries  References} % 参考文献加入到英文目录
\addtolength{\bibsep}{-0.8em}
\nocite{*}  %若将此命令屏蔽掉，则未引用的文献不会出现在文后的参考文献列表中。
%如果文档类是article之类的, 用\renewcommand\refname{参考文献}
%如果文档类是book之类的, 用\renewcommand\bibname{参考文献}
\bibliography{reference/reference}

% -*-coding: utf-8 -*-

\defaultfont
\appendix

\BiAppChapter{带章节的附录}{Full Appendix}%
完整的附录内容，包含章节，公式，图表等

%%%%%%%%%%%%%%%%%%%%%%%%%%%%%%%%%%%%%%%%%%%%%%%%%%%%%%%%%
\BiSection{附录节的内容}{Section in Appendix}
这是附录的节的内容

附录中公式的示例：
\begin{align}
a & = b \times c \\
E & = m c^2
\end{align}

% 致谢
%!Tex Program = xelatex
% -*-coding: utf-8 -*-
\defaultfont
\BiAppendixChapter{攻读硕士学位期间发表的论文及其他成果} {Papers
published in the period of PH.D. education}
\setlength{\parindent}{0em}
（一）发表的学术论文
\begin{publist}
\item 作者. 题目. 期刊. 年, 卷(期): 页码

\item 作者. 题目. 期刊. 年, 卷(期): 页码

\item 作者. 题目. 期刊. 年, 卷(期): 页码
\end{publist}

（二）申请及已获得的专利（无专利时此项不必列出）


（三）获得的科技奖励（无获奖时此项不必列出）

\vfill
\hangafter=1\hangindent=2em\noindent
注：发表的学术论文书写格式应该与参考文献相同，详细规范请参见~GB/T~7714-2005的顺序编码制形式；
如已发表的学术论文被~EI~或~SCI~收录，请标明收录号及~SCI~论文的影响因子；对已接收但尚未发表出来的学术论文，请注明是否~EI~或~SCI~刊源。


\setlength{\parindent}{2em}
    % 所发文章
%!Tex Program = xelatex
% -*-coding: utf-8 -*-
\phantomsection
\markboth{致谢}{致谢}
\addcontentsline{toc}{chapter}{\xiaosi 致谢}
\addcontentsline{toe}{chapter}{\bfseries \xiaosi Acknowledgements}  \chapter*{致谢}
\defaultfont



% 致谢
%!Tex Program = xelatex
% -*-coding: utf-8 -*-
\pagestyle{empty}
\newcommand{\subchapterstyle}%
  {\song\rmfamily\bfseries\xiaoer}  
\addcontentsline{toc}{chapter}{独创性声明}
\begin{center}{\subchapterstyle 独创性声明}\end{center}

 本人声明所呈交的学位论文是本人在导师指导下进行的研究工作及取得的研
 究成果。据我所知，除了文中特别加以标注和致谢的地方外，论文中不包含其他
 人已经发表或撰写过的研究成果，也不包含为获得\underline{\textbf{~东北林业大学~}}或其他教育
 机构的学位或证书而使用过的材料。与我一同工作的同志对本研究所做的任何贡献均己在论文中作了明确的说明并表示谢意

\vspace{\baselineskip}
\hspace{4em}学位论文作者签名：\hfill 签字日期：\hspace{2.5em}年\hspace{1.5em}月\hspace{1.5em}日

\vspace{2\baselineskip}
\addcontentsline{toc}{chapter}{学位论文版权使用授权书}
\begin{center}{\subchapterstyle 学位论文版权使用授权书}\end{center}

本学位论文作者完全了解\underline{\textbf{~东北林业大学~}}有关保留、使用学位论文的规
定，有权保留并向国家有关部门或机构送交论文的复印件和磁盘，允许论文被查
阅和借阅。本人授权\underline{\textbf{~东北林业大学~}}可以将学位论文的全部或部分内容编入有
关数据库进行检索，可以采用影印、缩印或扫描等复制手段保存、汇编学位论
文。
（保密的学位论文在解密后适用本授权书）

%\vspace{\baselineskip}
%本学位论文属于（请在以下相应方框内打“$\surd$”）：
%
%保密\xiaoer$\square$\xiaosi ，在~~~~~~~~~~~~~年解密后适用本授权书
%
%不保密\xiaoer$\square$\xiaosi

\vspace{\baselineskip}
\hspace{4em}学位论文作者签名：\hfill 签字日期：\hspace{2.5em}年\hspace{1.5em}月\hspace{1.5em}日

\vspace{\baselineskip}
\hspace{4em}导师签名：\hfill 签字日期：\hspace{2.5em}年\hspace{1.5em}月\hspace{1.5em}日

\vspace{2\baselineskip}
学位论文作者毕业后去向：\\

\vspace{\baselineskip}
\hspace{4em}工作单位：  \hfill 电话：\hspace{10.5em}

\vspace{\baselineskip}
\hspace{4em}通讯地址：  \hfill  邮编：\hspace{10.5em}

   % 承诺

\clearpage


\end{document}
